\documentclass[10pt]{article}
% Math Packages
\usepackage{amsmath, mathtools}
\usepackage{amssymb}
\usepackage{amsthm}
\usepackage{amsfonts}
\usepackage{bbm}
\usepackage{breqn}
\usepackage[margin=1in]{geometry}
\usepackage{graphicx}
\usepackage{tikz}
\usetikzlibrary{arrows.meta}
\usetikzlibrary{calc}
\usepackage{forest}
\usepackage{tikz-qtree}
\graphicspath{ {./images/} }
\usepackage{hyperref}
\usepackage[capitalize]{cleveref}
\usepackage[shortlabels]{enumitem}
\usetikzlibrary{arrows,matrix,positioning}
\usepackage{multicol}

% for the pipe symbol
\usepackage[T1]{fontenc}

% Citing theorems by name. (source: https://tex.stackexchange.com/questions/109843/cleveref-and-named-theorems)
\makeatletter
\newcommand{\ncref}[1]{\cref{#1}\mynameref{#1}{\csname r@#1\endcsname}}

\def\mynameref#1#2{%
  \begingroup
  \edef\@mytxt{#2}%
  \edef\@mytst{\expandafter\@thirdoffive\@mytxt}%
  \ifx\@mytst\empty\else
  \space(\nameref{#1})\fi
  \endgroup
}
\makeatother

% Colorful Notes
\usepackage{color} \definecolor{Red}{rgb}{1,0,0} \definecolor{Blue}{rgb}{0,0,1}
\definecolor{Purple}{rgb}{.5,0,.5} \def\red{\color{Red}} \def\blue{\color{Blue}}
\def\gray{\color{gray}} \def\purple{\color{Purple}}
\newcommand{\rnote}[1]{{\red [#1]}} % \rnote{foo} gives '[foo]' in red
\newcommand{\pnote}[1]{{\purple [#1]}} % \pnote{foo} gives '[foo]' in purple
\newcommand{\bnote}[1]{{\blue #1}} % \bnote{foo} gives 'foo' in blue
\newcommand{\gnote}[1]{{\gray #1}} % \gnote{foo} gives 'foo' in gray
\newcommand{\Max}[1]{{\purple [#1]}} % \bnote{foo} then 'foo' is blue


% Claim numbering (the counter restarts after each proof environment)
\newcounter{claimcount}
\setcounter{claimcount}{0}
\newenvironment{claim}{\refstepcounter{claimcount}\par\addvspace{\medskipamount}\noindent\textbf{Claim \arabic{claimcount}:}}{}
\usepackage{etoolbox}
\AtBeginEnvironment{proof}{\setcounter{claimcount}{0}}
\newenvironment{claimproof}{\par\addvspace{\medskipamount}\noindent\textit{Proof of Claim  \arabic{claimcount}.}}{\hfill\ensuremath{\qedsymbol} \tiny{Claim}

  \medskip}
% Add claim support to cleverref
\crefname{claimcount}{Claim}{Claims}


% Math Environments
\newtheorem{theorem}{Theorem}
\newtheorem{assumption}[theorem]{Assumption}
\newtheorem{lemma}[theorem]{Lemma}
\newtheorem{proposition}[theorem]{Proposition}
\newtheorem{corollary}[theorem]{Corollary}
\newtheorem{question}[theorem]{Question}
\theoremstyle{definition}
\newtheorem{definition}[theorem]{Definition}
\newtheorem{remark}[theorem]{Remark}
\newtheorem{example}[theorem]{Example}
\newtheorem{notation}[theorem]{Notation}
\newtheorem{problem}[theorem]{Problem}

% Matrices and Column Vectors. 
\usepackage{stackengine}
\setstackgap{L}{1.0\normalbaselineskip}
\usepackage{tabstackengine}
\setstackEOL{;}% row separator
\setstackTAB{,}% column separator
\setstacktabbedgap{1ex}% inter-column gap
\setstackgap{L}{1.5\normalbaselineskip}% inter-row baselineskip
\let\nmatrix\bracketMatrixstack  %Usage: \nmatrix{1,2,3\4,5,6}
\newcommand\cv[1]{\setstackEOL{,}\bracketMatrixstack{#1}} %usage: \cv{1,2,3}

% Custom Math Coqmmands
\newcommand{\vt}{\vskip 5mm} % vertical space
\newcommand{\fl}{\noindent\textbf} % first line
\newcommand{\Fl}{\vt\noindent\textbf} % first line with space above
\newcommand{\norm}[1]{\left\lVert#1\right\rVert} % norm
\newcommand{\pnorm}[1]{\left\lVert#1\right\rVert_p} % p-norm
\newcommand{\qnorm}[1]{\left\lVert#1\right\rVert_q} % q-norm
\newcommand{\1}[1]{\textbf{1}_{\left[#1\right]}} % indicator function 
\def\limn{\lim_{n\to\infty}} % shortcut for lim as n-> infinity
\def\sumn{\sum_{n=1}^{\infty}} % shortcut for sum from n=1 to infinity
\def\sumkn{\sum_{k=1}^{n}} % shortcut for sum from k=1 to n
\def\sumin{\sum_{i=1}^{n}} % shortcut for sum from i=1 to n
\def\SAs{\sigma\text{-algebras}} % shortcut for $\sigma$-algebras
\def\SA{\sigma\text{-algebra}} % shortcut for $\sigma$-algebra
\def\Ft{\mathcal{F}_t} % time-indexed sigma-algebra (t)
\def\Fs{\mathcal{F}_s} % time-indexed sigma-algebra (s)
\def\F{\mathcal{F}} % sigma-algebra
\def\G{\mathcal{G}} % sigma-algebra
\def\R{\mathbb{R}} % Real numbers
\def\N{\mathbb{N}} % Natural numbers
\def\Z{\mathbb{Z}} % Integers
\def\E{\mathbb{E}} % Expectation
\def\P{\mathbb{P}} % Probability
\def\Q{\mathbb{Q}} % Q probability
\def\dist{\text{dist}} %Text 'dist' for things like 'dist(x,y)'
\newcommand{\indep}{\perp \!\!\! \perp}  %independence symbol
\def\Var{\mathrm{Var}} % Variance
\def\tr{\mathrm{tr}} % trace

% Brackets and Parentheses
\def\[{\left [}
    \def\]{\right ]}
% \def\({\left (}
%   \def\){\right )}



\usepackage{color}
\definecolor{Red}{rgb}{1,0,0}
\definecolor{Blue}{rgb}{0,0,1}
\definecolor{Purple}{rgb}{.5,0,.5}
\def\red{\color{Red}}
\def\blue{\color{Blue}}
\def\gray{\color{gray}}
\def\purple{\color{Purple}}
\definecolor{RoyalBlue}{cmyk}{1, 0.50, 0, 0}
\newcommand{\dempfcolor}[1]{{\color{RoyalBlue}#1}} 
\newcommand{\demph}[1]{\dempfcolor{{\sl #1}}}

% comment exactly one of the following line to show / hide the solutions
% \newcommand{\solution}[1]{{\purple #1}} % uncomment to show the solutions
\newcommand{\solution}[1]{} % uncomment to hide the solutions



\title{Lecture Notes for Math 372: \\Elementary Probability and Statistics}
\date{Last updated: \today}
% \author{mh}

\begin{document}
\begin{center}
  \section*{Math 372: Homework 05}
  \textit{Due Friday, Feb 27}
\end{center}

\textit{Instructions: Do not write your solutions on this sheet---you don't have enough space to show
your work. So use separate sheets of paper.}



\begin{problem}
  The prevalance of a disease in a population is $0.15\%$ (i.e, $0.0015$). A new diagnostic test is
  advertised as having an false negative rate of $1\%$ and a false positive rate of $1\%$. What is the
  probability that someone who tests positive doesn't actually have the disease?

  % [Draw a tree diagram]
  
  % We want $\P\left[ \text{healthy}\mid\text{tests positive} \right]$. For brevity, let's let
  % \begin{equation*}
  %   H = \left[ \text{health} \right] \quad \text{and} \quad P = \left[ \text{tests positive} \right].
  % \end{equation*}
  % \begin{align}\label{eq:3}
  %   \P\left[ H\mid P\right]
  %   &= \frac{\P\left[ H\cap P \right]}{\P\left[P \right]}
  % \end{align}
  % Now,
  % \begin{equation*}
  %   \P\left[ H\cap P \right] = (0.9985)(0.01)= 0.009985
  % \end{equation*}
  % and by the Law of Total Probability, 
  % \begin{align*}
  %   \P\left[P \right] 
  %   &= \P\left[ H\cap P \right]
  %   + \P\left[ H^{c}\cap P \right] \\
  %   &= (0.9985)(0.01) + (0.0015)(0.99)
  % \end{align*}
  % Thus by \cref{eq:3},
  % \begin{equation*}
  %   \P\left[ H\mid P \right]
  %   = \frac{0.009985}{(0.9985)(0.01) + (0.0015)(0.99)} \approx 0.87.
  % \end{equation*}
  % Therefore, approximately 87\% of people who test positive aren't actually sick! 
\end{problem}


\begin{problem}[independence]
  Alice rolls six dice and wins if she rolls at least one six. Bob rolls
  twelve dice and wins if he rolls at least two sixes. Who has the greater
  probability of winning? \textit{[Hint: calculate the probabilities to lose.]}
\end{problem}


\begin{problem}
  \textit{(You do not need to show your work for this problem.)} Suppose we roll  5 six-sided dice.
  \begin{enumerate}[(a)]
    \item  Compute the following probabilities
    \begin{enumerate}[(i)]
      \item Three of the numbers are even
      \item At least 1 six is thrown
      \item No even numbers are thrown
    \end{enumerate}
    \item What is the expected number of sixes?
  \end{enumerate}
\end{problem}



\begin{problem}
  Compute $\E\left[X \right]$, $\E\left[X^{2} \right]$ and $\Var(X)$ for the
  following random variables:
  \begin{enumerate}[(a)]
    \item $X = (\text{value of one roll of a 6-sided dice})$.
    \item A fair coin toss.
    $Y = \left\{ \begin{array}{l@{\quad:\quad}l} 1 & \text{with probability
        }\frac{1}{2}\\0 & \text{with probability
        }\frac{1}{2} \end{array}\right.$
    \item A biased coin toss. $Z = \left\{ \begin{array}{l@{\quad:\quad}l} 1 &  \text{with
          probability }p\\0 & \text{with probability }1-p \end{array}\right.$
  \end{enumerate}
\end{problem}



\begin{problem}
  A certain brand of upright freezer is available in three different rated
  capacities: $16$,  $18$, and $20$ cubic feet. Let $X$ be the rated capacity
  of a freezer of this brand sold at a certain store. Suppose that $X$ has
  probability mass function
  \begin{center}
    \begin{tabular}[c]{c|ccc}
      $x$&16&18&20\\ 
      \hline
      $p(x)$&.2&.5&.3 
    \end{tabular}
  \end{center}
  \begin{enumerate}[(a)]
    \item Compute $\E\left[X \right], \E\left[X^{2} \right], \Var(X)$.
    \item If the price of a freezer having capacity $X$ is $70X-650$, what is
    the expected price paid by the next customer to buy a freezer?
    \item What is the variance of the price paid by the next customer?
    \item Suppose that although the rated capacity of a freezer is $X$, the
    actual capacity is $h(X)=X-.008X^{2}$. What is the expected actual
    capacity of the freezer purchased by the next customer.
  \end{enumerate}
\end{problem}

\newpage

\begin{problem}[Binomial random variable]
  Let $\alpha\in (0,1)$ be a fixed parameter. Suppose we flip a coin 4 times, each time with the probability
  of heads being $\alpha$. Let
  \begin{equation*}
    X = \text{ the number of heads out among the 4 flips.}
  \end{equation*}
  \begin{enumerate}[(a)]
    \item  Make a table like the following, including all 16 possible outcomes:
    \begin{table}[htb]
      \centering
        \begin{tabular}{|c|c|c|}
          \hline
          \textbf{Outcome} & \textbf{Value of }$X$ & \textbf{Probability} \\ \hline
          HHHH & 4 & $\alpha^{4} $ \\ \hline
          HHHT & 3 & $\alpha^{3}(1-\alpha) $ \\ \hline
          HHTH & 3 & $\alpha^{3}(1-\alpha) $ \\ \hline
          HHTT & 2 & $\alpha^{2}(1-\alpha)^{2} $ \\ \hline
          ... &... & ... \\ \hline
        \end{tabular}
    \end{table}
    
    \item From now on, let's assume that $\alpha=\frac{1}{2}$. In other words, we are flipping fair
    coins. Recalling that $p(x) = \P\left[X=x \right]$, fill out a table like the following:
    \begin{center}
      \begin{tabular}{|l|l|l|l|l|l|}
        \hline
        $x$    & 0 & 1 & 2 & 3 & 4 \\ \hline
        $p(x)$ &   &   &   &   &   \\ \hline
      \end{tabular}
    \end{center}
    \item Verify that $\sum_{x}p(x) =1.$
    \item Compute $\E\left[X \right]$ and $\E\left[X^{2} \right]$.
    \item Use your answer from the previous part to compute $\text{Var}(X)$.
  \end{enumerate} 
\end{problem}


\begin{problem}[Properties of the geometric distribution]
  Let $\alpha\in (0,1)$, and let $X$ be a random variable with probability mass function
  \begin{equation*}
    p(x) = \left\{ \begin{array}{l@{\quad}l} \left( 1-\alpha \right)^{x-1}\alpha  & \text{if }x=1,2,3\ldots\\0 & \text{otherwise.} \end{array}\right.
  \end{equation*}
  (A random variable with this probability mass function is called a \textit{geometric random variable} with success parameter $\alpha$.)

  \begin{enumerate}[(a)]
    \item We know that in general probability mass functions should always sum
    up to one. Let's check that it is indeed the case here. Specifically,
    carefully show that $p(1)+p(2)+p(3)+\ldots =1$, making sure to justify
    your steps.
    \item Next we will compute the expected value of $X$ using the formula 
    \begin{equation*}
      \E\left[X \right] := \sum_{x}x p(x),
    \end{equation*}
    where $x$ ranges over all possible values of $X$.
    \begin{enumerate}[(i)]
      \item \label{item:1} Expand out above summation, including the first five terms followed by ``+\dots''
      \item \label{item:2} The geometric series formula tells us that as long as $|x|<1$, the following
      formula holds:
      \begin{equation*}
        a + ax+ax^{2}+ax^{3}+ax^{4}+ax^{5}+\ldots = \frac{a}{1-x}.
      \end{equation*}
      Differentiate both sides of the above equation with respect to $x$. 
      \item Use your answers to parts \ref{item:1} and\ref{item:2} to deduce the value of
      $\E\left[X \right]$
      \item What is the value of $\E\left[X \right]$ if $\alpha=\frac{1}{2}$? What about if
      $\alpha=\frac{2}{3}$?
    \end{enumerate}
  \end{enumerate}
\end{problem}

% \begin{problem}
%   Consider rolling a six-sided dice. Come up with two events that are independent. 
% \end{problem}


% \begin{problem}
%   Let $X$ be a random variable, and let $\mu= \E\left[X \right]$. We know that
%   \begin{equation*}
%     \Var(X) = \E\left[(X-\mu)^{2} \right].
%   \end{equation*}
%   Prove that
%   \begin{equation*}
%     \Var(X) = \E\left[X^{2} \right] - \mu^{2}.
%   \end{equation*}
% \end{problem}


% \begin{problem}[Exploding Dice]
%   Compute the expectation of an exploding 6-sided dice.
%   \textit{\textbf{Note: }An `exploding' dice roll is a mechanic used in various tabletop games whereby the player rolls a dice,
%     and then if they rolled the maximum number on the die, they get to roll it again and add the new roll to
%     the old one. This can happen more than once.} 
% \end{problem}




\end{document}